% ===================================================================
% 磁性数据存储的物理原理：从微磁学到硬盘驱动器
% Physical Principles of Magnetic Data Storage:
%   From Micromagnetics to Hard Disk Drives
% ===================================================================

\documentclass[11pt,a4paper,twoside]{article}

% ---------- Encoding & Fonts ----------
\usepackage[UTF8]{ctex}
% CJK fonts auto-detected by ctex (Windows: SimSun/SimHei, macOS: Songti SC/Heiti SC)
\usepackage[T1]{fontenc}
\usepackage{amsthm}  % must precede newtxmath to avoid \openbox conflict
\usepackage{newtxtext,newtxmath}
\usepackage[scaled=0.92]{helvet}
\usepackage{courier}

% ---------- Layout ----------
\usepackage[margin=2.5cm,top=3cm,bottom=2.8cm]{geometry}
\usepackage{setspace}
\onehalfspacing
\usepackage{fancyhdr}
\pagestyle{fancy}
\fancyhf{}
\fancyhead[LE]{\small\itshape 第1篇\quad 磁性数据存储的物理原理}
\fancyhead[RO]{\small\itshape 从微磁学到硬盘驱动器}
\fancyfoot[CE,CO]{\thepage}
\renewcommand{\headrulewidth}{0.4pt}

% ---------- Math & Physics ----------
\usepackage{amsmath,amsfonts}
\usepackage{mathtools}
\usepackage{bm}
\usepackage{physics}
\usepackage{siunitx}
\sisetup{separate-uncertainty,per-mode=symbol}
\newcommand{\degC}{\ensuremath{{}^{\circ}\mathrm{C}}}
\DeclareSIUnit{\angstrom}{\text{\AA}}

% ---------- Graphics & Color ----------
\usepackage{graphicx}
\usepackage{xcolor}
\usepackage{tikz}
\usepackage{tikz-3dplot}
\usetikzlibrary{arrows.meta,shapes,decorations.pathmorphing,patterns,calc,positioning}
\usepackage{pgfplots}
\pgfplotsset{compat=1.18}
\usepackage{float}

% ---------- Tables ----------
\usepackage{array,booktabs,tabularx,multirow}
\usepackage{colortbl}

% ---------- Hyperlinks & Metadata ----------
\usepackage[hidelinks,colorlinks=true,linkcolor=blue!60!black,citecolor=green!50!black,urlcolor=blue!60!black]{hyperref}
\usepackage[capitalise,nameinlink,noabbrev]{cleveref}

% ---------- Custom envs ----------
% amsthm already loaded before newtxmath
\theoremstyle{definition}
\newtheorem{definition}{定义}[section]
\theoremstyle{plain}
\newtheorem{theorem}{定理}[section]
\newtheorem{lemma}[theorem]{引理}
\newtheorem{proposition}[theorem]{命题}
\theoremstyle{remark}
\newtheorem{remark}{注解}[section]
\usepackage{tcolorbox}
\tcbuselibrary{skins,breakable}
\newtcolorbox{physbox}[1][]{
  colback=blue!3!white, colframe=blue!50!black,
  fonttitle=\bfseries, title=物理注释, breakable, #1
}

% ---------- Bibliography ----------
\usepackage[backend=bibtex,style=ieee,sorting=none]{biblatex}
\addbibresource{refs-magnetic.bib}

% ---------- Title ----------
\title{
  {\Huge\bfseries 磁性数据存储的物理原理} \\[8pt]
  {\Large 从微磁学到硬盘驱动器} \\[12pt]
  {\normalsize Physical Principles of Magnetic Data Storage: \\
   From Micromagnetics to Hard Disk Drives}
}
\author{}
\date{\today}

\begin{document}
\maketitle

\begin{center}
\rule{0.8\textwidth}{0.5pt} \\[8pt]
{\small 凝聚态物理与信息存储交叉专题 · 第1篇} \\[4pt]
\end{center}

% ===================================================================
\begin{abstract}
\noindent
硬盘驱动器（HDD）是凝聚态物理在信息技术领域最成功的应用之一。本文从微磁学、量子输运和统计物理的
第一性原理出发，系统阐述磁性数据存储的完整物理链条。第一部分建立铁磁薄膜介质的微观理论：从
Heisenberg交换作用与Stoner巡游磁性出发，导出磁自由能泛函的五个基本项——交换能、磁晶各向异性能、
退磁能、Zeeman能和Dzyaloshinskii-Moriya相互作用能，进而建立Landau-Lifshitz-Gilbert（LLG）
方程作为描述磁化动力学的第一性原理方程。第二部分分析存储介质的超顺磁极限与Arrhenius-N\'{e}el
热激活模型，严格导出保持时间与晶粒体积的指数依赖关系，阐释磁记录三元悖论的物理根源。第三部分
深入写入过程的Stoner-Wohlfarth相干翻转模型，给出astroid临界曲线的完整推导，并讨论热辅助磁
记录（HAMR）对超顺磁极限的突破——包括近场光学换能器的等离激元局域和FePt-$L1_0$有序相的超高
各向异性起源。第四部分详述隧道磁阻（TMR）读头：从Julli\`{e}re态密度模型到MgO(001)相干隧穿
的Bloch态对称性筛选理论，给出发射光谱与复能带结构的定量分析。第五部分讨论记录介质的空间组织、
伺服定位系统和信号处理链。本文力图在\"四大力学加固体物理\"的深度上，为磁性存储提供一个自洽
而完备的理论框架。

\medskip
\noindent\textbf{关键词}：微磁学；Landau-Lifshitz-Gilbert方程；Stoner-Wohlfarth模型；
超顺磁性；隧道磁阻；垂直磁记录；热辅助磁记录
\end{abstract}

\tableofcontents
\newpage

% ===================================================================
\section{引言}
\label{sec:intro}

从物理学的视角审视硬盘驱动器，它本质上是一个\textbf{可控的、可读出的磁化状态存储系统}。
信息被编码在铁磁薄膜中纳米级区域的磁化方向上。要完整理解这一系统，至少需要调用：
电磁学（写磁场的产生与聚焦）、量子力学（交换作用与自旋极化隧穿）、统计物理（热激活翻转与
Arrhenius-N\'{e}el动力学）、固体物理（能带结构与Bloch态对称性）、以及非线性动力学
（磁化进动与Gilbert阻尼）。

本文的叙述路径严格遵循物理因果链：
\[
\text{介质微观物理（磁化稳定性的量子起源）}
\to \text{写入物理（可控翻转的动力学）}
\to \text{读取物理（无损探测的输运理论）}
\to \text{系统集成（空间组织与信号处理）}
\]
每个环节都力图从第一性原理出发，推导出工程上可用的定量结论。各节之间的逻辑依赖如\cref{fig:logic}所示。

\begin{figure}[H]
\centering
\begin{tikzpicture}[
  node distance=2.2cm,
  box/.style={rectangle,draw=blue!60,fill=blue!5,rounded corners,text width=3.2cm,
              align=center,minimum height=1.2cm,font=\small\bfseries},
  arrow/.style={->,>=Stealth,line width=1pt,blue!60}
]
\node[box] (A) {海森堡交换作用 \\ + Stoner模型 \\[2pt] \S 2.1--2.3};
\node[box,right=of A] (B) {磁自由能泛函 \\ + LLG方程 \\[2pt] \S 2.3--3.1};
\node[box,below right=of B] (C) {Stoner-Wohlfarth翻转 \\ + 超顺磁极限 \\[2pt] \S 3.2--3.3};
\node[box,below left=of B] (D) {写磁头电磁学 \\ + HAMR等离激元 \\[2pt] \S 3.4--3.5};
\node[box,below=of C,xshift=-2cm] (E) {TMR读头 \\ Bloch态对称筛选 \\[2pt] \S 4};
\node[box,below=of E] (F) {介质噪声 \\ + 伺服定位 \\[2pt] \S 5--6};
\draw[arrow] (A) -- (B);
\draw[arrow] (B) -- (C);
\draw[arrow] (B) -- (D);
\draw[arrow] (C) -- (E);
\draw[arrow] (D) -- (E);
\draw[arrow] (E) -- (F);
\end{tikzpicture}
\caption{本文各节的物理逻辑依赖关系。箭头表示\"物理上依赖\"，非严格数学推导顺序。}
\label{fig:logic}
\end{figure}

% ===================================================================
\section{铁磁薄膜介质的微观物理基础}
\label{sec:microphysics}

\subsection{铁磁性的量子力学起源：交换作用与Stoner巡游磁性}

\subsubsection{Heisenberg局域矩图像}

铁磁性的本质是多体量子效应。在绝缘磁体（如YIG、铁氧体）中，磁性来源于局域d/f电子，
Heisenberg模型提供了精确的出发点。相邻自旋间的各向同性交换Hamiltonian为：
\begin{equation}\label{eq:heisenberg}
  \mathcal{H}_{\text{ex}} = -J \sum_{\langle i,j\rangle} \mathbf{S}_i \cdot \mathbf{S}_j
  = -J \sum_{\langle i,j\rangle} \left[
    S_i^z S_j^z + \frac{1}{2}\left(S_i^+ S_j^- + S_i^- S_j^+\right)
  \right]
\end{equation}
其中 $J > 0$ 为铁磁交换积分，求和限于最近邻。对于Co、Fe等3d过渡金属，
$J \sim \SI{10}{-20}{meV}$（由局域自旋密度泛函计算得到）。

交换作用的物理根源是Coulomb排斥与Pauli不相容原理的协同效应。考虑两个氢原子——Heitler-London
模型——单重态（$S=0$, 空间对称）与三重态（$S=1$, 空间反对称）的能量差为：
\begin{equation}
  J_{\text{HL}} = \frac{1}{2}(E_S - E_T)
  = \frac{1}{1 - S_{ab}^2}\left[ K_{ab} - S_{ab} J_{ab} \right]
\end{equation}
其中 $S_{ab} = \langle a|b \rangle$ 为交叠积分，$K_{ab} = \langle ab|1/r_{12}|ba \rangle$
为交换积分，$J_{ab} = \langle ab|1/r_{12}|ab \rangle$ 为Coulomb积分。铁磁性要求
$J_{\text{HL}} > 0$，即交换积分压倒交叠项——这一条件在3d金属中因d轨道的空间局域性和
Fermi面附近态密度高峰而自然满足。

\begin{physbox}
  \textbf{铁磁性的竞争判据：Bethe-Slater曲线。}
  交换积分 $J$ 的值和符号依赖于原子间距 $R$ 与d轨道半径 $r_d$ 的比值 $R/r_d$。
  当 $R/r_d > 3$（Co, Ni）时 $J > 0$（铁磁耦合）；当 $R/r_d < 3$（Mn, Cr）时
  $J < 0$（反铁磁耦合）。Co的 $R/r_d \approx 3.3$，恰好落在铁磁区。
  这一经验规律解释了为何CoFe和NiFe合金是优良的软磁材料，
  而纯Mn是反铁磁的——尽管它们都是3d过渡金属。
\end{physbox}

\subsubsection{Stoner巡游电子模型}

对于金属铁磁体（Fe, Co, Ni及其合金），磁性电子并非局域而是巡游的——它们占据扩展的Bloch态。
Stoner模型（1938）提供了巡游铁磁性的平均场图像。在密度泛函理论（DFT）的现代视角下，
Stoner模型相当于自旋极化的局域密度近似（LSDA）的最简形式。

设自旋向上（$\uparrow$）和自旋向下（$\downarrow$）的态密度分别为
$D_\uparrow(E)$ 和 $D_\downarrow(E)$。在无自旋极化时两自旋通道的态密度相等：
$D_\uparrow(E) = D_\downarrow(E) = D(E)/2$。交换劈裂 $\Delta E_{\text{ex}}$ 使两自旋能带
发生刚性相对位移，产生自旋不平衡的占据数：
\begin{align}
  n_\uparrow &= \int_{-\infty}^{E_F} D\left(E - \frac{\Delta E_{\text{ex}}}{2}\right)\, dE \\
  n_\downarrow &= \int_{-\infty}^{E_F} D\left(E + \frac{\Delta E_{\text{ex}}}{2}\right)\, dE
\end{align}
净磁矩 $m = n_\uparrow - n_\downarrow$（以 $\mu_B$ 为单位）。自洽性要求交换劈裂
$\Delta E_{\text{ex}} = I m$，其中 $I$ 为Stoner交换参数。对于小 $m$ 展开，
自洽条件化为：
\begin{equation}\label{eq:stoner-criterion}
  I \cdot D(E_F) > 1
\end{equation}
此即\textbf{Stoner判据}。对于bcc Fe，$D(E_F) \approx \SI{1.5}{eV^{-1}}$，
$I \approx \SI{0.9}{eV}$，故 $I D(E_F) \approx 1.35 > 1$，满足巡游铁磁性条件。
对于Co（hcp或fcc），由于 $E_F$ 位于d带的高态密度区域，
$I D(E_F) \approx 1.6$——更强的铁磁倾向，解释了Co基合金中更大的 $M_s$ 和 $K_u$。

每Co原子的磁矩在Slater-Pauling曲线上的值约为 \SI{1.72}{\mu_B}，这与实验值
\SI{1.67}{\mu_B}（hcp Co）吻合良好。注意Slater-Pauling曲线的物理：在强铁磁体中，
多数自旋d带完全填满（每个原子约5个d电子），剩余的少数自旋d电子数与平均价电子数的关系
线性递减。Co恰好在曲线峰值附近——这是其广泛用于磁记录介质的根本电子结构原因。

\subsection{磁自由能泛函与连续介质微磁学}

\subsubsection{从离散自旋到连续场：微磁学近似}

在微磁学中（Brown, 1963），我们采用连续介质描述，将磁化矢量视为位置的光滑函数
$\mathbf{M}(\mathbf{r})$。这一近似的合法性依赖于\textbf{交换长度}
$l_{\text{ex}}$ 远大于原子间距的条件。单位磁化矢量的方向场定义为：
\begin{equation}
  \mathbf{m}(\mathbf{r}) = \frac{\mathbf{M}(\mathbf{r})}{M_s}, \quad |\mathbf{m}(\mathbf{r})| = 1
\end{equation}
其中 $M_s = |\mathbf{M}|$ 为饱和磁化强度（对于CoCrPt-氧化物复合薄膜，
$M_s \approx \SI{3e5}{A/m}$，具体值取决于Pt含量和氧化物掺量）。

\subsubsection{自由能泛函的五项分解}

总磁Gibbs自由能泛函 $E_{\text{total}}[\mathbf{m}]$ 可分解为五个物理来源不同的贡献：
\begin{equation}\label{eq:free-energy}
  E_{\text{total}} = E_{\text{ex}} + E_{\text{ani}} + E_{\text{demag}}
                   + E_{\text{Zeeman}} + E_{\text{DMI}}
\end{equation}
以下逐一详述。

\paragraph{(a) 交换能 $E_{\text{ex}}$}

在连续极限下，$\mathbf{m}$ 的空间变化导致交换能的升高（类似于非线性 $\sigma$ 模型的梯度项）：
\begin{equation}\label{eq:exchange}
  E_{\text{ex}} = A \int_V (\nabla \mathbf{m})^2 \, d^3r
  = A \int_V \left[ (\nabla m_x)^2 + (\nabla m_y)^2 + (\nabla m_z)^2 \right] d^3r
\end{equation}
其中 $A$ 为交换刚度常数（对于hcp Co，$A \approx \SI{3.0e-11}{J/m}$；
对于CoCrPt合金，$A \approx \SI{1.5e-11}{J/m}$，因非磁性氧化物掺入稀释了交换耦合）。
$A$ 与Heisenberg交换积分 $J$ 的关系为（对于简单立方晶格，晶格常数 $a$）：
$A \approx J S^2 z / (2a)$，其中 $z$ 为配位数。

上述梯度项使 $\nabla \mathbf{m}$ 的非均匀性付出能量代价，因此交换能倾向于使 $\mathbf{m}$ 在空间中
均匀一致——这是铁磁有序在连续介质层面的体现。但交换作用本身是各向同性的，不能单独确定磁化方向。

\paragraph{(b) 磁晶各向异性能 $E_{\text{ani}}$}

在垂直磁记录介质（CoCrPt合金）中，hcp结构的c轴（$[0001]$方向）优先沿垂直于膜面取向，
这是由外延生长和薄膜应力共同诱导的织构。对于单轴各向异性——垂直磁记录的主导项：
\begin{equation}\label{eq:anisotropy}
  E_{\text{ani}} = K_u \int_V \sin^2\theta(\mathbf{r}) \, d^3r
\end{equation}
其中 $\theta$ 为 $\mathbf{m}$ 与易磁化轴（垂直方向）的夹角，
$K_u$ 为单轴各向异性常数（单位 J/m$^3$）。

$K_u$ 的微观物理起源包含三个机制：

\begin{enumerate}
  \item \textbf{磁晶各向异性}：自旋-轨道耦合（$\mathcal{H}_{\text{SO}} = \xi \mathbf{L}\cdot\mathbf{S}$）
  将各向同性的自旋系统\"钉扎\"到晶格的方向上。当磁化方向旋转时，电子云通过自旋-轨道耦合重新
  调整轨道角动量，从而改变轨道-晶格场（晶体电场，CEF）的耦合能。
  对于hcp Co，体磁晶各向异性常数 $K_1^{\text{bulk}} \approx \SI{4.5e5}{J/m^3}$。

  \item \textbf{界面（N\'{e}el）各向异性}：在多层膜中（如Co/Pt或Co/Pd），
  界面处Co原子的配位环境不对称（一侧为Co，另一侧为Pt/Pd），
  导致界面Co的d轨道择优取向。N\'{e}el（1954）提出了基于原子对方向有序的各向异性模型：
  界面各向异性能 $K_s$ 正比于界面处Co-Pt对的键各向异性。
  对于Co/Pt界面，$K_s \approx \SI{1}{-2}{mJ/m^2}$。
  当Co层厚度减小到 \SI{<1}{nm} 时，界面贡献可以超越体贡献，
  使得有效 $K_u^{\text{eff}} = K_u^{\text{bulk}} + 2K_s/t_{\text{Co}}$
  显著增大——此即界面各向异性工程。

  \item \textbf{磁弹各向异性}：薄膜生长过程中的残余应变通过磁致伸缩耦合到磁化方向。
  对于磁致伸缩常数 $\lambda_s$ 的材料，各向异性增量
  $\Delta K_{\text{me}} = (3/2) \lambda_s \sigma$，
  其中 $\sigma$ 为面内应力。CoCrPt的 $\lambda_s \approx \SI{-2e-5}{}$（负值），
  因此在面内压应力下垂直各向异性增强。
\end{enumerate}

对于现代垂直记录介质（CoCrPt-SiO$_2$），$K_u \approx \SIrange{2}{5e5}{J/m^3}$。
而HAMR介质（$L1_0$有序FePt）的 $K_u \approx \SI{7e6}{J/m^3}$——超过CoCrPt的十倍——
其超大各向异性的根源在于Fe-3d与Pt-5d的强自旋-轨道杂化，
详见\cref{sec:hamr}。

\paragraph{(c) 退磁能 $E_{\text{demag}}$}

退磁能（又称静磁能或形状各向异性能）是长程偶极-偶极相互作用的连续极限。
磁化体内的等效磁荷密度为 $\rho_m = -\nabla \cdot \mathbf{M}$，
面磁荷密度为 $\sigma_m = \mathbf{M} \cdot \hat{n}$。退磁场 $\mathbf{H}_d$ 由磁标势
$\phi_M$ 的梯度给出：
\begin{align}
  \mathbf{H}_d &= -\nabla \phi_M \label{eq:hd} \\
  \nabla^2 \phi_M &= \nabla \cdot \mathbf{M} \label{eq:poisson-mag}
\end{align}
退磁能密度为：
\begin{equation}\label{eq:demag}
  \varepsilon_{\text{demag}} = -\frac{\mu_0}{2} \mathbf{M} \cdot \mathbf{H}_d
\end{equation}
因子 $1/2$ 避免了双重计数（自能修正）。

对于垂直磁化的无限大薄板（厚度 $t$），退磁场在板内是均匀的：
$\mathbf{H}_d = -M_s \hat{z}$，退磁能密度为 $\mu_0 M_s^2/2$。
对于CoCrPt，$M_s \approx \SI{3.0e5}{A/m}$，
$\mu_0 M_s^2/2 \approx \SI{5.7e4}{J/m^3}$。

正是退磁能的这一正值使垂直于膜面的磁化处于不利的能量状态——
形状各向异性倾向于将磁化拉回面内。
垂直记录之所以可行，是因为磁晶各向异性 $K_u$ 压倒退磁能：
\begin{equation}
  K_{\text{eff}} = K_u - \frac{\mu_0}{2} M_s^2 > 0
  \label{eq:keff}
\end{equation}
这是垂直磁记录介质设计的第一个判据。对于现代介质，$K_{\text{eff}} \approx \SIrange{1}{4e5}{J/m^3}$，
对应的有效各向异性场 $H_K^{\text{eff}} = 2 K_{\text{eff}} / (\mu_0 M_s) \approx \SI{8}{-20}{kOe}$。

\paragraph{(d) Zeeman能 $E_{\text{Zeeman}}$}

外磁场 $\mathbf{H}_{\text{ext}}$（写磁场）与磁化的耦合能——写入操作的驱动源：
\begin{equation}
  E_{\text{Zeeman}} = -\mu_0 \int_V \mathbf{M}(\mathbf{r}) \cdot \mathbf{H}_{\text{ext}}(\mathbf{r}) \, d^3r
\end{equation}

\paragraph{(e) Dzyaloshinskii-Moriya相互作用 $E_{\text{DMI}}$}

在缺乏反演对称性的体系中（如重金属/铁磁界面），反对称交换作用以如下泛函形式出现：
\begin{equation}
  E_{\text{DMI}} = D \int_V \left[ m_z (\nabla \cdot \mathbf{m})
                   - (\mathbf{m} \cdot \nabla) m_z \right] d^3r
\end{equation}
其中 $D$ 为DMI常数（J/m$^2$）。DMI倾向于使相邻自旋倾斜（产生N\'{e}el型畴壁手性），
是磁斯格明子的来源。在传统HDD记录介质中，DMI是非主导项（$D \ll \sqrt{A K_u}$），
但在未来的赛道存储器中有重要作用。此处引用仅为完整性。

\subsubsection{有效场与Euler-Lagrange方程}

由自由能泛函对 $\mathbf{M}$ 取泛函导数（并保持 $|\mathbf{m}| = 1$ 的约束），
得到有效场的表达式：
\begin{equation}
  \mathbf{H}_{\text{eff}}(\mathbf{r}) = -\frac{1}{\mu_0}\frac{\delta E_{\text{total}}}{\delta \mathbf{M}(\mathbf{r})}
\end{equation}
代入五个贡献后：
\begin{equation}\label{eq:heff-full}
  \boxed{
  \mathbf{H}_{\text{eff}} = \mathbf{H}_{\text{ext}} + \mathbf{H}_d
  + \frac{2A}{\mu_0 M_s^2} \nabla^2 \mathbf{m}
  + \frac{2K_u}{\mu_0 M_s^2} (\mathbf{m} \cdot \hat{u}) \hat{u}
  + \frac{D}{\mu_0 M_s} \left[ (\nabla \cdot \mathbf{m})\hat{z} - \nabla m_z \right]
  }
\end{equation}
其中 $\hat{u}$ 为单轴各向异性的易磁化方向单位矢量（垂直记录中 $\hat{u} = \hat{z}$），
第三项（交换场）来自泛函导数中的分部积分（$\delta(\nabla\mathbf{m})^2/\delta\mathbf{m} = -2\nabla^2\mathbf{m}$）。

\subsection{磁滞回线：存储的物理基础}

铁磁体区别于顺磁体的本质特征是\textbf{磁滞}：$M(H)$ 是多值的，
系统的当前状态依赖于磁场的历史——这正是\"记忆\"的物理定义。
一个理想化（单畴）粒子的磁滞回线如\cref{fig:hysteresis}所示。

\begin{figure}[H]
\centering
\begin{tikzpicture}[scale=1.2]
  \draw[->] (-4,0) -- (4,0) node[right] {$H$};
  \draw[->] (0,-2.2) -- (0,2.5) node[above] {$M / M_s$};
  \draw[thick, red!70!black]
    (-3.5,-1.5) -- (1.5,-1.5) -- (1.5,1.5) -- (-1.5,1.5) -- (-1.5,-1.5) -- (-3.5,-1.5);
  \draw[dashed, blue!50] (-3.5,-1.5) -- (0,-1.5) -- (0,1.5) -- (2.5,1.5);
  \node[red!70!black] at (3.5,1.5) {矩形回线};
  \node[blue!50] at (3.5,0.8) {实际回线};
  \draw[<->] (1.5,0.2) -- (0,0.2) node[midway,above] {$H_c$};
  \node at (0.3,1.6) {$+M_r$};
  \node at (0.3,-1.6) {$-M_r$};
\end{tikzpicture}
\caption{理想矩形磁滞回线（红色）与典型实际回线（蓝色虚线）的对比。
$M_r$ 为剩余磁化强度，$H_c$ 为矫顽力。矩形比 $S = M_r/M_s$ 越接近1，
两个稳定态的可区分性越好。}
\label{fig:hysteresis}
\end{figure}

磁滞回线的关键参数及其物理意义：

\begin{itemize}
  \item \textbf{饱和磁化强度 $M_s$}：所有原子磁矩平行排列时的最大磁化强度，
  由材料成分决定（Slater-Pauling曲线）。在绝对零度下，
  $M_s(0) = n_{\text{atom}} \mu_{\text{atom}}$。
  \item \textbf{剩余磁化强度 $M_r$}：外场归零后的可保持磁化强度——信息存储的直接物理量。
  $M_r < M_s$（非理想矩形比）源于晶粒间的磁化方向分散。
  \item \textbf{矫顽力 $H_c$}：使 $M = 0$ 所需的反向磁场。
  度量数据抵抗外界干扰（杂散磁场、热涨落）的能力。
  由\cref{sec:sw}的Stoner-Wohlfarth模型给出微观表达式。
  \item \textbf{矩形比 $S = M_r / M_s$}：当 $S \to 1$ 时两个稳定态（$\pm M_r$）的
  磁场窗口最大，读信噪比最优。
\end{itemize}

% ===================================================================
\section{磁化动力学与写入过程}
\label{sec:writing}

\subsection{Landau-Lifshitz-Gilbert（LLG）方程}
\label{sec:llg}

磁化矢量 $\mathbf{M}(\mathbf{r}, t)$ 的时间演化由LLG方程描述——这是微磁学的基础动力学
方程，也是所有磁化翻转模拟的出发点：
\begin{equation}\label{eq:llg}
  \boxed{
  \frac{\partial \mathbf{M}}{\partial t}
  = -\gamma \mathbf{M} \times \mathbf{H}_{\text{eff}}
  + \frac{\alpha}{M_s} \mathbf{M} \times \frac{\partial \mathbf{M}}{\partial t}
  }
\end{equation}
其中：
\begin{itemize}
  \item $\gamma = g \mu_B / \hbar \approx \SI{1.7609e11}{rad.s^{-1}.T^{-1}}$ 为
  自由电子的旋磁比。对于Co，Land\'{e}因子 $g \approx 2.17$（轨道角动量未完全猝灭），
  $\gamma_{\text{Co}} \approx \SI{1.91e11}{rad.s^{-1}.T^{-1}}$。
  \item $\mathbf{H}_{\text{eff}}$ 为有效场，由\cref{eq:heff-full}给出。
  \item $\alpha$ 为Gilbert阻尼常数（无量纲），对于CoCrPt，
  $\alpha \approx 0.02$--$0.05$，由铁磁共振（FMR）线宽测量确定。
\end{itemize}

LLG方程的物理图像：
\begin{itemize}
  \item \textbf{进动项} $-\gamma \mathbf{M} \times \mathbf{H}_{\text{eff}}$：
  描述磁化绕有效场的无耗散Larmor进动。在 $\mathbf{H}_{\text{eff}} = \SI{1}{T}$
  的场中，进动角频率 $\omega = \gamma H_{\text{eff}} \approx \SI{1.76e11}{rad/s}$
  （即 \SI{28}{GHz}——微波范围）。
  \item \textbf{阻尼项} $(\alpha/M_s) \mathbf{M} \times \partial \mathbf{M}/\partial t$：
  源自自旋-晶格弛豫——进动能量通过自旋-轨道耦合传递给声子（晶格振动），
  最终以热的形式耗散。阻尼使 $\mathbf{M}$ 逐渐趋向 $\mathbf{H}_{\text{eff}}$ 的方向，
  即自由能的局部极小点。阻尼的微观机制包括：
  \begin{itemize}
    \item 自旋-轨道耦合引起的磁弹弛豫（Kittel机制）。
    \item s-d交换相互作用引起的传导电子自旋翻转散射（Mitchell机制）。
    \item 磁子-声子散射。
  \end{itemize}
\end{itemize}

在数值微磁学中（如使用OOMMF或Mumax3），LLG方程通常在有限差分或有限元网格上离散，
使用Runge-Kutta-Fehlberg（RKF45）或Heun格式进行时间积分。
典型的时间步长 $\Delta t \sim \SI{0.1}{-1}{ps}$（受进动频率限制：
$\Delta t < 1/(\gamma H_{\text{max}})$ 以保证数值稳定）。

\subsubsection{LLG的Landau-Lifshitz等价形式}

在某些文献中，LLG方程被重写为Landau-Lifshitz（LL）形式：
\begin{equation}
  \frac{\partial \mathbf{M}}{\partial t}
  = -\frac{\gamma}{1+\alpha^2} \mathbf{M} \times \mathbf{H}_{\text{eff}}
  - \frac{\gamma\alpha}{(1+\alpha^2) M_s}
    \mathbf{M} \times (\mathbf{M} \times \mathbf{H}_{\text{eff}})
\end{equation}
LL形式的优点在于阻尼项显式地表达为驱使 $\mathbf{M}$ 趋向 $\mathbf{H}_{\text{eff}}$
方向的力矩（$\mathbf{M} \times (\mathbf{M} \times \mathbf{H}_{\text{eff}})$
是 $\mathbf{H}_{\text{eff}}$ 垂直于 $\mathbf{M}$ 的分量，方向指向 $\mathbf{H}_{\text{eff}}$），
物理图像更为直观。两种形式通过 $1+\alpha^2$ 重标度等价——在 $\alpha \ll 1$
（实际材料的常见情形）下近似相同。

\subsection{Stoner-Wohlfarth相干翻转模型}
\label{sec:sw}

\subsubsection{单畴粒子假设}

对于晶粒尺寸 $d < l_{\text{ex}}$（交换长度，\SI{5}{-10}{nm}）的纳米粒子，
交换能远大于退磁能的梯度，所有原子自旋被迫同步旋转——即\textbf{相干翻转}。
Stoner-Wohlfarth（SW）模型（1948年，Phil. Trans. R. Soc. A）在此假设下
可解析求解，揭示了单畴粒子磁化翻转的几何基础。

\subsubsection{能量与平衡条件}

考虑一个单轴各向异性的椭球单畴粒子。令易磁化轴方向为 $\hat{u}$（假定沿 $\hat{z}$），
磁化方向为 $\mathbf{m} = (\sin\theta\cos\phi, \sin\theta\sin\phi, \cos\theta)$，
外加均匀磁场 $\mathbf{H}$ 与易磁化轴的夹角为 $\psi$。
总磁能密度为：
\begin{equation}\label{eq:sw-energy}
  \varepsilon(\theta) = K_u \sin^2\theta - \mu_0 M_s H \cos(\theta - \psi)
\end{equation}
其中第一项为各向异性能（两个极小值在 $\theta = 0, \pi$），
第二项为Zeeman能。

平衡态 $\partial\varepsilon/\partial\theta = 0$ 给出：
\begin{equation}\label{eq:sw-equilibrium}
  2K_u \sin\theta \cos\theta + \mu_0 M_s H \sin(\theta - \psi) = 0
\end{equation}
稳定性条件 $\partial^2\varepsilon/\partial\theta^2 > 0$：
\begin{equation}\label{eq:sw-stability}
  2K_u (\cos^2\theta - \sin^2\theta) + \mu_0 M_s H \cos(\theta - \psi) > 0
\end{equation}

引入无量纲场 $\mathbf{h} = \mathbf{H} / H_K$，其中各向异性场
$H_K = 2K_u / (\mu_0 M_s)$。\cref{eq:sw-equilibrium} 化为：
\begin{equation}
  \sin\theta \cos\theta + h_x \sin\theta \sin\phi - h_z \cos\theta = 0
\end{equation}
（取 $\mathbf{h} = h_x \hat{x} + h_z \hat{z}$，即外场在 $xz$ 平面内）。

\subsubsection{Astroid临界曲线}

翻转发生的条件为稳定性边界——\cref{eq:sw-stability}的等号成立时。
在 $h_x$-$h_z$ 平面上，该边界构成一条四尖点星形线（astroid）：
\begin{equation}\label{eq:astroid}
  \boxed{ h_x^{2/3} + h_z^{2/3} = 1 }
\end{equation}
当外场矢量 $\mathbf{H}$ 的箭头落在astroid外部时，只有一个稳定平衡态（翻转已经完成）；
落在内部时存在两个稳定平衡态（取决于历史）。

astroid的几何特征（\cref{fig:astroid}）：
\begin{itemize}
  \item 在 $\mathbf{h} = (0, -1)$ 处（即场沿易磁化轴反方向时），矫顽力取最大
  $H_c^{\text{max}} = H_K$。
  \item 在场方向与易轴成 \ang{45} 时，矫顽力最小：
  $H_c^{\text{min}} = H_K / 2$。
  \item astroid在 $|h_x| = 1$ 处的尖点对应于 $H = H_K$ 的垂直翻转——
  当场恰垂直于易轴时，稳定性以连续方式过渡，无滞后。
\end{itemize}

\begin{figure}[H]
\centering
\begin{tikzpicture}[scale=4]
  \draw[->] (-1.3,0) -- (1.3,0) node[right] {$h_x$};
  \draw[->] (0,-1.3) -- (0,1.3) node[above] {$h_z$};
  % Astroid
  \draw[thick, blue!70!black, domain=0:360, samples=200, variable=\t]
    plot ({cos(\t)^3}, {sin(\t)^3});
  % Cusp points
  \filldraw (1,0) circle (0.4pt) node[below right] {$(1,0)$};
  \filldraw (0,1) circle (0.4pt) node[above right] {$(0,1)$};
  \filldraw (-1,0) circle (0.4pt) node[below left] {$(-1,0)$};
  \filldraw (0,-1) circle (0.4pt) node[below left] {$H_c^{\text{max}}$};
  % H_K
  \draw[->,red!70!black,thick] (0,0) -- (0,-1);
  \node[red!70!black] at (0.2,-0.5) {$H_K$};
  % Stable vs unstable
  \node[blue!70!black] at (-0.5,-0.5) {2个稳定态};
  \node[blue!70!black] at (-1.0,0.8) {1个稳定态};
\end{tikzpicture}
\caption{Stoner-Wohlfarth astroid（星形线）。曲线内部区域（蓝色）对应两个稳定磁化方向，
曲线外部对应唯一稳定方向。磁化翻转发生在场矢量 $\mathbf{H}$ 穿越astroid时。}
\label{fig:astroid}
\end{figure}

\subsubsection{对写磁头设计的启示}

写磁头必须产生足够强的磁场，使在astroid外部的场分量足以翻转沿任意方向排列的晶粒。
对于沿易磁化轴反方向的晶粒（$\psi = \pi$，最\"顽固\"的取向），
所需写场 $H_{\text{write}} > H_K$。对于随机取向的晶粒系综，
最坏情况所需的写场介于 $H_K/2$ 和 $H_K$ 之间。
在介质设计中 $H_K \approx \SI{15}{-20}{kOe}$，
因此写磁头极尖磁场需达到 \SI{15}{-20}{kOe}——这是写磁头材料（需高 $M_s$ 和
高磁导率的软磁材料，如NiFe、CoFe、FeCoN等）设计的关键约束。

\subsection{热激活翻转与Arrhenius-N\'{e}el定律}
\label{sec:superpara}

\subsubsection{超顺磁极限的导出}

在绝对零度下，$K_u V \gg k_B T$ 保证了磁化方向被各向异性能垒牢牢限制在两个易磁化方向之一。
但在室温下，热涨落可以以概率 $\propto \exp(-K_u V / k_B T)$ 克服能垒。
N\'{e}el（1949，Ann. G\'{e}ophys.）建立了热激活翻转的Arrhenius定律：
\begin{equation}\label{eq:neel}
  \boxed{ \tau = \tau_0 \exp\!\left( \frac{K_u V}{k_B T} \right) }
\end{equation}
其中：
\begin{itemize}
  \item $\tau$ 为平均翻转时间（即在无外场条件下磁化自发反转的特征等待时间）。
  \item $\tau_0$ 为尝试时间（attempt time），与进动频率相关：
  $\tau_0^{-1} \sim \gamma H_K / (2\pi) \sim \SI{e9}{-e10}{s^{-1}}$，
  故 $\tau_0 \approx \SIrange{e-10}{e-9}{s}$。
\end{itemize}

以 $\tau > \SI{10}{years} \approx 3.15 \times 10^8\ \si{s}$ 为数据保持要求，
得到稳定性判据：
\begin{equation}\label{eq:stability-criterion}
  \boxed{ \frac{K_u V}{k_B T} \gtrsim 60 }
\end{equation}
（注：按上述 $\tau_0$ 与保持时间，单晶粒下限为 $\ln(\tau/\tau_0) \approx 40$；
``60''为工业界计入晶粒体积分布与可靠性裕度后的工程判据。）

该判据直接揭示了\textbf{磁记录的三元悖论}（magnetic recording trilemma）：

\begin{enumerate}
  \item \textbf{密度提升} $\rightarrow$ 更小晶粒体积 $V$（每个比特由有限个晶粒平均）。
  \item \textbf{热稳定性} $\rightarrow$ 更小的 $V$ 使能垒 $K_u V$ 下降，
  必须增大 $K_u$ 以维持 $K_u V / k_B T > 60$。
  \item \textbf{可写性} $\rightarrow$ 更大的 $K_u$ 导致更高的矫顽力
  $H_c \propto K_u / M_s$，需要更大的写磁场——但写磁头极尖磁场受软磁材料的
  $M_s$（约 \SI{2.4}{T} 为 FeCo 的极限）约束。
\end{enumerate}

三元悖论是推动HAMR和MAMR技术的根本物理动机。

\subsubsection{热场辅助翻转：N\'{e}el-Brown公式}

当外场 $\mathbf{H}$ 施加在易磁化轴反方向时，能垒被外场降低。
在SW近似下，降垒因子为 $(1 - H/H_K)^n$，其中指数 $n$ 取决于场的角度和SW对称性。
修正的N\'{e}el-Brown翻转时间为：
\begin{equation}\label{eq:neel-brown}
  \tau(H) = \tau_0 \exp\!\left[ \frac{K_u V}{k_B T}
            \left( 1 - \frac{H}{H_K} \right)^n \right]
\end{equation}
对于场严格按照易轴反方向，$n = 2$（因为能垒在SW模型中随 $H$ 以 $(1 - H/H_K)^2$ 标度）；
对于随机取向的系综平均，有效 $n \approx 1.5$（Victora, 1989, PRL）。

有限温度下的翻转不再是确定性的，而是\textbf{随机过程}。
翻转概率 $P_{\text{sw}}(t_p)$ 在脉冲宽度 $t_p$ 内为：
\begin{equation}
  P_{\text{sw}}(t_p) = 1 - \exp\!\left[ -\frac{t_p}{\tau(H)} \right]
\end{equation}
这意味着即使 $H < H_K$（在SW零温极限下不会翻转），
仍存在非零的翻转概率——热涨落在\"帮助\"系统越过能垒。
此即\textbf{热辅助写入}的根本物理：可以在低于 $H_K$ 的写场下，依靠有限温度实现充分高
的翻转概率。但这一过程也会引起热稳定性问题——相邻磁道上的不需要的翻转（相邻道擦除，ATE）。

\begin{physbox}
  \textbf{热稳定性的数值感觉。}在 $T=\SI{300}{K}$，$k_B T = \SI{4.14e-21}{J}
  \approx \SI{25.8}{meV}$。为达到 $\tau > \SI{10}{yr}$（即 $K_u V / k_B T > 60$），
  需要 $K_u V > \SI{2.48e-19}{J}$。对于CoCrPt晶粒（$K_u = \SI{3e5}{J/m^3}$），
  最小晶粒直径 $d_{\text{min}} = (6V/\pi)^{1/3} \approx \SI{8.5}{nm}$。
  当前 \SI{1}{Tb/in^2} 面密度的比特单元含约10--15个晶粒，每个晶粒直径约
  \SI{8}{-10}{nm}——已逼近超顺磁极限。再提升面密度必然需要更高的 $K_u$ 介质
  （如FePt, $K_u = \SI{7e6}{J/m^3}$）配合HAMR写入。
\end{physbox}

\subsection{写磁头的电磁学}
\label{sec:write-head}

写磁头是一个微米尺度的电磁铁。在准静态近似下（磁头尺寸远小于电磁波长，
$\lambda = c/f \sim \SI{30}{cm}$ 对应 \SI{1}{GHz}，远大于磁头尺寸
$\sim\SI{10}{\mu m}$），写场由静磁学Maxwell方程描述：
\begin{align}
  \nabla \times \mathbf{H} &= \mathbf{J} \label{eq:ampere} \\
  \nabla \cdot \mathbf{B} &= 0 \label{eq:gauss} \\
  \mathbf{B} &= \mu_0 (\mathbf{H} + \mathbf{M}_{\text{pole}}) \label{eq:constitutive}
\end{align}
其中 $\mathbf{J}$ 为线圈电流密度，$\mathbf{M}_{\text{pole}}$ 为磁极材料的磁化强度。

\subsubsection{Karlqvist场的解析近似}

在简化的二维缝隙模型中（Karlqvist, 1954），垂直磁记录写磁头的极尖区域被视为一个宽度为 $g$
（write gap）的缝隙，两壁的磁标势分别设为 $\pm V_0/2$。缝隙正下方深度为 $z$ 处的垂直场分量为：
\begin{equation}\label{eq:karlqvist}
  H_z(x,z) = \frac{H_0}{\pi} \left[
    \arctan\!\left(\frac{x + g/2}{z}\right)
    - \arctan\!\left(\frac{x - g/2}{z}\right)
  \right]
\end{equation}
其中 $H_0 = V_0 / g$ 为缝隙中心面处的场强（Amp\`{e}re定律估算：
$H_0 \approx N I / g$，$N$ 为线圈匝数，$I$ 为写电流）。
该式的物理图像：在间隙正下方场最强，沿横向半宽约 $g/2$ 后衰减。

对于现代写磁头，$g \approx \SI{30}{-50}{nm}$，飞行高度（磁头到介质面的间距）
$z \approx \SI{3}{-5}{nm}$，因此磁头产生的场在介质面有极高的梯度
（$dH_z/dx \sim H_0 / g$），这对写入陡峭的比特边界至关重要。

\subsubsection{软磁底层（SUL）的镜像效应}

现代垂直记录采用单极型写磁头配合盘片下方的软磁底层（SUL, Soft UnderLayer）。
SUL是高磁导率（$\mu_r \gg 1$）的软磁薄膜，其作用可用镜像法理解：
单极子在SUL表面产生一个符号相同的镜像单极子，等效于写极尖的有效长度加倍，
从而大幅增强了记录层中的写场强度和垂直均匀性。SUL的厚度需大于磁通回路的截面积要求
（典型 \SI{30}{-80}{nm}），材料为NiFe或CoFe非晶合金。

\subsection{热辅助磁记录（HAMR）}
\label{sec:hamr}

\subsubsection{基本原理：瞬态降低矫顽力}

HAMR的核心思想优雅而直接：写入时局部加热记录介质以瞬态降低 $H_c$，
使中等强度的写场足以翻转高 $K_u$ 介质；冷却后 $H_c$ 恢复，数据稳定保持。
加热温度 $T_{\text{write}}$ 需满足：
\begin{equation}
  H_c(T_{\text{write}}) < H_{\text{write}} < H_c(T_{\text{ambient}})
\end{equation}
其中 $T_{\text{write}}$ 通常在 \SI{400}{-700}{\degC{}}，
取决于介质材料的 $H_c(T)$ 曲线（\cref{fig:hc-vs-t}）。

\begin{figure}[H]
\centering
\begin{tikzpicture}
  \draw[->] (0,0) -- (8,0) node[right] {$T$ [\si{\degC{}}]};
  \draw[->] (0,0) -- (0,5) node[above] {$H_c$ [T]};
  \draw[thick, blue!70!black]
    (0.5,4.5) -- (2,4.3) -- (4,3.5) -- (5.5,2.0) -- (7,0.3);
  \draw[dashed, red!70!black] (0.2,1.8) -- (7.8,1.8);
  \node[red!70!black,right] at (7.8,1.8) {$H_{\text{write}}$};
  \draw[<->, thin] (3,0.3) -- (3,4.3)
    node[midway,right,font=\footnotesize] {可写};
  \draw[<->, thin] (0.5,0.3) -- (0.5,4.5)
    node[midway,left,font=\footnotesize] {稳定};
  \node[blue!70!black] at (7,0.6) {$T_C$};
  \draw[thin,dotted] (5.5,0) -- (5.5,2.0);
  \node at (5.5,-0.3) {$T_{\text{write}}$};
\end{tikzpicture}
\caption{矫顽力随温度的变化曲线（蓝色）。在HAMR写入温度 $T_{\text{write}}$ 附近，
$H_c$ 降至写磁头能力（红色虚线）以下，写入可行。
冷却至环境温度后，$H_c$ 恢复至远高于写场，数据稳定。}
\label{fig:hc-vs-t}
\end{figure}

\subsubsection{近场光学换能器的等离激元物理}

在 \SI{405}{-830}{nm} 的写入激光波长下，衍射极限光斑直径约为 $\lambda/(2\text{NA}) \approx
\SI{200}{-400}{nm}$，远大于现代磁记录的比特尺寸（\SI{20}{-30}{nm}）。
HAMR通过\textbf{近场光学换能器（NFT, Near-Field Transducer）}将光斑压缩至亚衍射极限。

NFT是集成在写磁头滑块上的等离子体纳米天线（通常为Au或Ag）。激光聚焦到NFT上后，
激发表面等离激元共振——自由电子的集体振荡。在NFT的尖端（曲率半径 \SI{5}{-10}{nm}），
局域表面等离激元共振（LSPR）产生极强的近场增强：
\begin{equation}
  |\mathbf{E}_{\text{local}}|^2 \sim \eta_{\text{NFT}} \times |\mathbf{E}_{\text{incident}}|^2
\end{equation}
增强因子 $\eta_{\text{NFT}}$ 可达 $10^3$--$10^4$，使得光斑的有效加热区域缩小至
\SI{20}{-30}{nm}——即亚波长的热点。

加热媒介的吸收功率密度为：
\begin{equation}
  Q(\mathbf{r}) = \frac{1}{2} \omega \varepsilon_0 \operatorname{Im}[\varepsilon(\omega)]
                 |\mathbf{E}_{\text{local}}(\mathbf{r})|^2
\end{equation}
由于金属（Au、Ag以及介质中的FePt颗粒）在近红外至可见光波段有有限的
$\operatorname{Im}[\varepsilon]$，光能通过欧姆损耗转化为热。

\subsubsection{热传导与淬冷动力学}

介质内的温度时空演化由热传导方程描述：
\begin{equation}\label{eq:heat}
  \rho C_p \frac{\partial T}{\partial t}
  = \nabla \cdot (\boldsymbol{\kappa} \cdot \nabla T) + Q(\mathbf{r}, t)
\end{equation}
其中 $\rho C_p$ 为体比热容，$\boldsymbol{\kappa}$ 为热导张量
（面内热导远大于面外热导——由薄膜的热边界和界面散射引起，
$\kappa_{\parallel} / \kappa_{\perp} \approx 3$--$5$）。

关键热时间常数由热扩散长度 $\ell_{\text{th}} = \sqrt{\kappa \tau / (\rho C_p)}$ 估算。
对于 \SI{10}{nm} 厚的记录层，面外散热时间 $\sim \SI{100}{ps}$——这意味着HAMR的热脉冲
必须精确同步于写场脉冲，时间抖动需控制在 $\pm\SI{50}{ps}$ 以内。
淬冷速率 $dT/dt \sim \SI{e10}{-e11}{K/s}$，远高于非晶化所需的临界冷却速率，
但足够快以\"冻结\"目标磁化方向。

\subsubsection{FePt $L1_0$ 有序相的物理}

HAMR的使能材料是\textbf{$L1_0$ 有序FePt合金}。$L1_0$ 相为四方晶系
（空间群 $P4/mmm$，$a = \SI{3.85}{\angstrom}$, $c = \SI{3.71}{\angstrom}$），
Fe和Pt原子沿c轴交替排列。

$L1_0$ FePt的极高 $K_u \approx \SI{7e6}{J/m^3}$ 的微观起源是双重的：
\begin{enumerate}
  \item \textbf{晶体电场效应}：Fe-3d能级在四方晶场中劈裂为 $d_{xy}$、
  $d_{xz/yz}$、$d_{x^2-y^2}$ 和 $d_{3z^2-r^2}$ 子能级。
  四方畸变使 $d_{3z^2-r^2}$ 和 $d_{x^2-y^2}$ 的能级分离，
  产生了单轴各向异性的电子结构基础。
  \item \textbf{Fe-3d/Pt-5d杂化}：Pt的5d态具有巨大的自旋-轨道耦合常数
  ($\xi_{\text{Pt}} \approx \SI{0.5}{eV}$，远高于Fe-3d的 \SI{0.05}{eV})。
  Fe-3d与Pt-5d的强烈杂化将Pt的SOC\"借入\"Fe的磁性电子，
  大幅增强了各向异性能。第一性原理LSDA计算（Daalderop et al., 1991, PRB；
  Sakuma, 1994, JPSJ）确认 $K_u$ 的主要贡献来自少数自旋的Fe $d_{xy}/d_{x^2-y^2}$
  与Pt $d_{xy}$ 的杂化态对SOC的响应。
\end{enumerate}

$L1_0$ FePt的居里温度 $T_C \approx \SI{750}{K}$（\SI{477}{\degC{}}）。
在HAMR操作中，介质被局部加热至 $T_{\text{write}} \approx \SI{600}{-700}{K}$，
此时 $M_s$ 下降至室温值的 $40\%$--$20\%$，$H_c$ 从室温的 \SI{70}{kOe}
降至 \SI{5}{-10}{kOe}——落入写头的磁场能力范围。
写入后急速冷却（通过散热层热容，\cref{eq:heat}），介质恢复至平行于外场的磁化方向。

HAMR技术的商业化（Seagate Exos Mozaic 3+, 2024上市）标志着磁记录密度突破
\SI{3}{Tb/in^2} 的开始，是近场光学与微磁学深度交叉的里程碑。

% ===================================================================
\section{读取过程：自旋极化输运与隧道磁阻}
\label{sec:reading}

\subsection{巨磁阻（GMR）效应的物理基础}

GMR效应（A. Fert 和 P. Gr\"{u}berg, 1988, PRL——2007年诺贝尔物理学奖）
是自旋电子学的奠基石。其物理核心是\textbf{自旋相关散射}。

\subsubsection{Mott的双电流模型}

Mott（1936）提出：在低温下（$T \ll T_C$），自旋翻转散射的速率远低于动量散射速率，
因此自旋向上和自旋向下的传导电子可被视为两个独立的并联导电通道：
\begin{equation}
  \sigma = \sigma_\uparrow + \sigma_\downarrow
\end{equation}
在铁磁金属中，自旋向上和自旋向下的电阻率通常不同（$\rho_\uparrow \neq \rho_\downarrow$），
这是GMR效应的根源。

\subsubsection{GMR自旋阀的电阻网络}

GMR自旋阀的基本结构：铁磁层1 (F1) / 非磁间隔层 (NM, Cu) / 铁磁层2 (F2)。
当F1和F2的磁化平行（P态）时，多数自旋电子在两层中均属多数自旋，
经历低散射 $\rightarrow$ 低电阻 $R_P$。当反平行（AP态）时，每种自旋至少在一层中
为少数自旋（经历强散射） $\rightarrow$ 高电阻 $R_{AP}$。

在宏观二流体电阻网络模型中：
\begin{equation}
  R_P = \frac{2R_\uparrow R_\downarrow}{R_\uparrow + R_\downarrow}, \quad
  R_{AP} = \frac{R_\uparrow + R_\downarrow}{2}
\end{equation}
（假设两铁磁层厚度相同，界面散射主导）。
GMR比值：
\begin{equation}
  \frac{\Delta R}{R_P} = \frac{R_{AP} - R_P}{R_P}
  = \frac{(R_\uparrow - R_\downarrow)^2}{4 R_\uparrow R_\downarrow}
\end{equation}
典型自旋阀中 $\Delta R/R_P \approx 10\%$--$20\%$（电流垂直于膜面, CPP配置）。

\subsection{隧道磁阻（TMR）与MgO(001)相干隧穿}
\label{sec:tmr}

\subsubsection{Julli\`{e}re模型（1975）}

在F1 / 绝缘势垒(I) / F2结构的磁隧道结（MTJ）中，Julli\`{e}re提出了基于态密度的简单模型。
假定隧穿电导正比于发射极和收集极的态密度之积：
\begin{equation}
  G_P \propto D_{1\uparrow} D_{2\uparrow} + D_{1\downarrow} D_{2\downarrow}, \quad
  G_{AP} \propto D_{1\uparrow} D_{2\downarrow} + D_{1\downarrow} D_{2\uparrow}
\end{equation}
定义电极的自旋极化率：
\begin{equation}
  P_i = \frac{D_{i\uparrow}(E_F) - D_{i\downarrow}(E_F)}
              {D_{i\uparrow}(E_F) + D_{i\downarrow}(E_F)}
\end{equation}
则TMR比值为：
\begin{equation}\label{eq:julliere}
  \boxed{ \text{TMR} = \frac{R_{AP} - R_P}{R_P}
  = \frac{2 P_1 P_2}{1 - P_1 P_2} }
\end{equation}
对于传统AlO$_x$非晶势垒 + 3d铁磁金属（Co, Fe, NiFe），
$P \approx 0.4$--$0.5$（TMR $\approx 50\%$--$70\%$）。

\subsubsection{MgO(001)的Bloch态对称性过滤理论}

2004年，Parkin（IBM）和Yuasa（AIST）独立报道了
Fe(001)/MgO(001)/Fe(001)和CoFeB/MgO/CoFeB外延/织构结的室温TMR > 200\%，
打破了Julli\`{e}re模型的预测上限。这一突破的根源是\textbf{Bloch态对称性筛选效应}，
由Butler等人（2001, PRB）和Mathon与Umerski（2001, PRB）
的第一性原理Layer-KKR计算预言。

物理图景（\cref{fig:tmr-bloch}）：

\begin{enumerate}
  \item MgO的复能带结构在其带隙内具有倏逝态。在这些倏逝态中，
  \textbf{$\Delta_1$ 对称性}（s-like，$k_\parallel = 0$ 点的全对称不可约表示）
  的衰减常数最小（$\kappa_{\Delta_1} \approx \SI{0.2}{\per\angstrom}$），
  因此经MgO势垒传输的距离最长。

  \item 体心立方（bcc）Fe的多数自旋能带中，在 $E_F$ 处存在 $\Delta_1$ 态的Bloch波函数——
  主要由 $s$、$p_z$、$d_{3z^2-r^2}$ 杂化构成。少数自旋能带中，$E_F$ 处的Bloch态
  主要为 $\Delta_2$ 和 $\Delta_5$ 对称性（由 $d_{xz/yz}$ 等轨道的杂化构成），
  没有 $\Delta_1$ 分量。

  \item 因此：
  \begin{itemize}
    \item \textbf{平行态（P）}：两侧Fe均为多数自旋 $\rightarrow$
    $\Delta_1 \to \Delta_1 \to \Delta_1$，高隧穿概率（$\propto e^{-\kappa_{\Delta_1} d}$）。
    \item \textbf{反平行态（AP）}：一侧为多数自旋 $\Delta_1$，另一侧为少数自旋
    （$\Delta_2, \Delta_5$）$\rightarrow$ 对称性不匹配，隧穿被指数级压制。
  \end{itemize}
\end{enumerate}

定量上，隧穿电导（Landauer-B\"{u}ttiker公式在散射矩阵表述中）为：
\begin{equation}
  G = \frac{e^2}{h} \sum_{\mathbf{k}_\parallel} T(\mathbf{k}_\parallel)
\end{equation}
其中 $T(\mathbf{k}_\parallel)$ 为横向波矢 $\mathbf{k}_\parallel$
的态的总传输概率，由左右电极Bloch态与势垒倏逝态的匹配边界条件决定
（通过Layer-KKR Green函数方法或第一性原理量子输运程序计算）。

\begin{figure}[H]
\centering
\begin{tikzpicture}[scale=1.0]
  % Left electrode
  \fill[blue!20] (0,0) rectangle (1.5,4);
  \node at (0.75,2) {Fe (多数自旋)};
  \node at (0.75,3.5) {$\Delta_1$};
  % MgO barrier
  \fill[yellow!30] (1.5,0) rectangle (3.0,4);
  \node at (2.25,2) {MgO};
  \node at (2.25,3.5) {$\Delta_1$ (低衰减)};
  % Right electrode
  \fill[blue!20] (3.0,0) rectangle (4.5,4);
  \node at (3.75,2) {Fe (多数自旋)};
  \node at (3.75,3.5) {$\Delta_1$};
  % Arrows
  \draw[->,very thick,green!60!black] (1.5,3.5) -- (3.0,3.5);
  \node at (2.25,3.8) {\textbf{平行态: 高透射}};

  % AP case below
  \fill[red!20] (0,-1.5) rectangle (1.5,0.5);
  \node at (0.75,-0.5) {Fe (多数自旋)};
  \node at (0.75,0.2) {$\Delta_1$};
  \fill[yellow!30] (1.5,-1.5) rectangle (3.0,0.5);
  \node at (2.25,-0.5) {MgO};
  \fill[red!10] (3.0,-1.5) rectangle (4.5,0.5);
  \node at (3.75,-0.5) {Fe (少数自旋)};
  \node at (3.75,-0.1) {$\Delta_2, \Delta_5$};
  \draw[->,very thick,red!70!black] (1.5,0.2) -- (2.5,0.2);
  \draw[red!70!black,->] (2.5,0.2) -- (2.5,-1.5) node[below] {衰减};
  \node at (2.25,0.5) {\textbf{反平行态: 低透射}};
\end{tikzpicture}
\caption{TMR的Bloch态对称性过滤机制的示意图。
在平行态（上），$\Delta_1 \to \Delta_1 \to \Delta_1$ 路径畅通；
在反平行态（下），由于少数自旋能带中缺少 $\Delta_1$ 态，对称性不匹配导致隧穿概率急降。}
\label{fig:tmr-bloch}
\end{figure}

\subsubsection{MgO-TMR的数值表现}

当前最优的CoFeB/MgO/CoFeB MTJ的室温TMR可超过 $600\%$（对应
$\Delta R/R_P > 6$），电阻面积积（RA product）约为
1--10~$\Omega\cdot\mu$m$^2$。RA值与MgO厚度近似指数关系
$RA \propto \exp(\kappa d)$，典型 $d_{\text{MgO}} \approx \SI{1.0}{-1.5}{nm}$。

\subsection{读磁头的信号生成与噪声}

TMR读头的自由层磁化受盘片介质杂散场驱动而旋转。
读出电压信号由TMR效应转换：
\begin{equation}
  \Delta V = V_{\text{bias}} \times \frac{\Delta R}{R}
           \approx V_{\text{bias}} \times \frac{\text{TMR}}{2} (1 - \cos\theta)
\end{equation}
其中 $\theta$ 为自由层与参考层磁化的相对夹角。
小偏转角（线性响应区）下，$\theta \propto H_s$（杂散场），
因此 $\Delta V \propto H_s$。典型信号幅度约 \SI{1}{-5}{mV}。

读回信噪比（SNR）的噪声源层次分析：

\begin{enumerate}
  \item \textbf{介质噪声}（主导）：晶粒尺寸和位置的无规分布导致读出信号的涨落。
  介质SNR正比于 $\sqrt{N_{\text{grains}}}$（$N_{\text{grains}}$ 为每比特内晶粒数）。
  随密度提升，$N_{\text{grains}}$ 下降，介质噪声恶化——这是记录密度最基本的信噪比瓶颈。
  \item \textbf{Johnson热噪声}（Nyquist定理）：
  $\langle V_n^2 \rangle = 4 k_B T R \Delta f$。
  对于 $R \approx \SI{100}{\Omega}$，$\Delta f \approx \SI{1}{GHz}$
  （\SI{2}{Gb/s} 数据率），$V_n \approx \SI{40}{\mu V}$（均方根）。
  \item \textbf{散粒噪声}（Shot noise）：$\langle I_n^2 \rangle = 2 e I \Delta f$。
  对于 $I \approx \SI{1}{mA}$（偏置电流），$I_n \approx \SI{0.6}{\mu A}$，
  等效电压噪声 $V_n = I_n R \approx \SI{60}{\mu V}$——与热噪声（\SI{40}{\mu V}）量级相当。
  \item \textbf{磁噪声}（mag-noise）：自由层磁化的热激发进动。
  在FMR频率附近有频谱峰，是GHz频段的重要噪声源。
\end{enumerate}

\subsection{从原始BER到用户BER：LDPC纠错编码}

原始读出信号的位错误率（raw BER）约为 $10^{-3}$--$10^{-2}$（受介质噪声限制）。
现代HDD使用\textbf{低密度奇偶校验码（LDPC）}——一种稀疏图码——通过迭代置信度传播
（Belief Propagation）解码，将用户可见BER降至 $< 10^{-15}$。

LDPC解码在物理层面利用了来自TMR读头的\textbf{软信息}（对数似然比, LLR）：
\begin{equation}
  \text{LLR}_i = \ln \frac{P(y_i | x_i = 0)}{P(y_i | x_i = 1)}
\end{equation}
其中 $y_i$ 为读回电压采样值，$x_i$ 为写入比特的真实值。
LLR的绝对值量度了该比特判决的\"可靠程度\"，
软判决LDPC比硬判决Viterbi解码获得约 \SI{2}{-3}{dB} 的编码增益。
这是\"信道编码定理\"（Shannon, 1948）在磁记录系统中的最成功实现。

% ===================================================================
\section{数据组织与物理寻址}
\label{sec:organization}

\subsection{磁道与扇区的空间布局}

盘片表面被精密划分为一组同心圆磁道，每条磁道被周向分割为扇区。
现代硬盘道间距约 \SI{50}{-70}{nm}（对应TPI约 \num{350000}--\num{500000}），
外圈磁道的周向扇区数多于内圈——因为外圈半径更大、线速度更高、
信号质量更好（Zone Bit Recording, ZBR）。

\subsection{伺服定位的物理学}

为了将磁头定位至宽度仅数十纳米的磁道上空，需要高带宽伺服控制系统。
位置误差信号（PES）通过读取伺服帧中预写入的burst图案来生成：
\begin{equation}
  \text{PES} = \frac{V_A - V_B}{V_A + V_B}
\end{equation}
其中 $V_A, V_B$ 分别是伺服burst A和B（写入在目标磁道两侧半道距处）的读回幅度。
PES经PID控制器放大后驱动音圈电机（VCM）和微动器。

音圈电机的驱动力为Lorentz力：$\mathbf{F} = N I \mathbf{L} \times \mathbf{B}_{\text{PM}}$。
对于道密度 \SI{500}{ktpi}（道间距 \SI{50}{nm}），
定位精度要求 $3\sigma \approx \SI{\pm5}{nm}$——这需要亚纳米级的机械稳定性和伺服带宽
\SI{2}{-5}{kHz}。实际系统中，可重复偏离（RRO）和不可重复偏离（NRRO）的物理来源
包括盘片颤动、轴承噪声、空气湍流和外部振动。

% ===================================================================
\section{结论与展望}
\label{sec:conclusion}

本文所建立的磁性存储物理框架可总结为以下核心链条：

\begin{enumerate}
  \item \textbf{存储介质}：铁磁性来源于3d电子的交换作用（Stoner判据
  $I D(E_F) > 1$）；纳米晶粒通过磁晶各向异性 $K_u$ 和形状各向异性定义了
  两个剩磁态（$\pm M_r$）；热稳定性由Arrhenius-N\'{e}el定律
  $\tau = \tau_0 \exp(K_u V / k_B T)$ 保证，$K_u V/k_B T > 60$ 对应
  $\tau > \SI{10}{yr}$。

  \item \textbf{写入}：LLG方程 $\partial_t \mathbf{M} = -\gamma \mathbf{M}\times\mathbf{H}_{\text{eff}}
  + (\alpha/M_s)\mathbf{M}\times\partial_t \mathbf{M}$ 统一描述进动与阻尼；
  Stoner-Wohlfarth astroid $h_x^{2/3} + h_z^{2/3} = 1$ 给出了翻转的临界条件；
  HAMR通过加热瞬态降低 $H_c$ 以突破超顺磁极限。

  \item \textbf{读取}：TMR结中MgO(001)势垒充当Bloch态对称性过滤器——
  $\Delta_1$ 态的低衰减常数和高透射选择性在平行态产生低电阻、
  反平行态产生高电阻；TMR > 600\% 的室温值使信噪比足以支持
  \SI{>1}{Tb/in^2} 的可靠检测。

  \item \textbf{动力学}：LLG方程是微磁学模拟的第一性原理方程，
  结合热噪声项（满足涨落-耗散定理）的随机LLG方程可定量模拟
  写入错误率和热稳定性。

  \item \textbf{极限}：磁记录的三元悖论（密度-稳定性-可写性）驱动着
  HAMR、BPM（比特图案化介质）和MAMR（微波辅助磁记录）等新一代技术的物理创新。
\end{enumerate}

\begin{physbox}
  \textbf{与后续论文的衔接。}本文建立了磁性存储的完整物理图景。
  理解这些原理是分析\"删除\"操作为何不能立即消除盘片磁化、
  以及剩磁可以被MFM探测的物理前提。
  详见第5篇《数据删除与恢复的物理原理》。
\end{physbox}

% ===================================================================
\appendix
\section{关键物理常数与材料参数}

\begin{table}[H]
\centering
\caption{磁性存储涉及的关键物理参数}
\begin{tabular}{@{}llll@{}}
\toprule
参数 & 符号 & 典型值 & 单位 \\
\midrule
Bohr磁子 & $\mu_B$ & $9.274 \times 10^{-24}$ & J/T \\
旋磁比（自由电子） & $\gamma_e$ & $1.761 \times 10^{11}$ & rad/(s$\cdot$T) \\
Co饱和磁化强度 & $M_s$(Co) & $1.42 \times 10^6$ & A/m \\
CoCrPt 各向异性能密度 & $K_u$ & $(2\text{--}5)\times 10^5$ & J/m$^3$ \\
FePt($L1_0$) 各向异性能密度 & $K_u$ & $7 \times 10^6$ & J/m$^3$ \\
Co 交换刚度 & $A$ & $3.0 \times 10^{-11}$ & J/m \\
Gilbert阻尼常数（CoCrPt） & $\alpha$ & 0.02--0.05 & — \\
MgO势垒高度（对Co） & $\phi_B$ & 3.2--3.7 & eV \\
MgO $\Delta_1$ 态衰减常数 & $\kappa_{\Delta_1}$ & $\approx 0.2$ & \AA$^{-1}$ \\
当前晶粒尺寸 & $d$ & 6--10 & nm \\
当前道间距 & — & 50--70 & nm \\
当前面密度（PMR） & — & 1.0--1.5 & Tb/in$^2$ \\
\bottomrule
\end{tabular}
\end{table}

% ===================================================================
\begin{thebibliography}{99}

\bibitem{aharoni}
A. Aharoni, \textit{Introduction to the Theory of Ferromagnetism}, 2nd ed.
Oxford University Press, 2000.
\href{https://scholar.google.com/scholar?q=Introduction+to+the+Theory+of+Ferromagnetism+Aharoni}{Google Scholar}

\bibitem{bertram}
H. N. Bertram, \textit{Theory of Magnetic Recording}.
Cambridge University Press, 1994.
\href{https://doi.org/10.1017/CBO9780511623066}{doi:10.1017/CBO9780511623066}

\bibitem{brown}
W. F. Brown, Jr., \textit{Micromagnetics}.
Interscience Publishers, 1963.
\href{https://scholar.google.com/scholar?q=Micromagnetics+Brown+1963}{Google Scholar}

\bibitem{stoner48}
E. C. Stoner and E. P. Wohlfarth,
``A mechanism of magnetic hysteresis in heterogeneous alloys,''
\textit{Phil. Trans. R. Soc. Lond. A}, vol. 240, pp. 599--642, 1948.
\href{https://doi.org/10.1098/rsta.1948.0007}{doi:10.1098/rsta.1948.0007}

\bibitem{neel49}
L. N\'{e}el,
``Th\'{e}orie du tra\^{i}nage magn\'{e}tique des ferromagn\'{e}tiques en grains fins avec applications aux terres cuites,''
\textit{Ann. G\'{e}ophys.}, vol. 5, pp. 99--136, 1949.
\href{https://scholar.google.com/scholar?q=Neel+trainage+magnetique+ferromagnetiques+grains+fins+1949}{Google Scholar}

\bibitem{slonczewski96}
J. C. Slonczewski,
``Current-driven excitation of magnetic multilayers,''
\textit{J. Magn. Magn. Mater.}, vol. 159, pp. L1--L7, 1996.
\href{https://doi.org/10.1016/0304-8853(96)00062-5}{doi:10.1016/0304-8853(96)00062-5}

\bibitem{butler01}
W. H. Butler, X.-G. Zhang, T. C. Schulthess, and J. M. MacLaren,
``Spin-dependent tunneling conductance of Fe$|$MgO$|$Fe sandwiches,''
\textit{Phys. Rev. B}, vol. 63, 054416, 2001.
\href{https://doi.org/10.1103/PhysRevB.63.054416}{doi:10.1103/PhysRevB.63.054416}

\bibitem{parkin04}
S. S. P. Parkin, C. Kaiser, A. Panchula, P. M. Rice, B. Hughes, M. Samant, and S.-H. Yang,
``Giant tunnelling magnetoresistance at room temperature with MgO (100) tunnel barriers,''
\textit{Nature Materials}, vol. 3, pp. 862--867, 2004.
\href{https://doi.org/10.1038/nmat1256}{doi:10.1038/nmat1256}

\bibitem{kryder08}
M. H. Kryder, E. C. Gage, T. W. McDaniel, W. A. Challener, R. E. Rottmayer, G. Ju, Y.-T. Hsia, and M. F. Erden,
``Heat assisted magnetic recording,''
\textit{Proc. IEEE}, vol. 96, pp. 1810--1835, 2008.
\href{https://doi.org/10.1109/JPROC.2008.2004315}{doi:10.1109/JPROC.2008.2004315}

\bibitem{richards59}
B. Richards and E. Wolf,
``Electromagnetic diffraction in optical systems II. Structure of the image field in an aplanatic system,''
\textit{Proc. Roy. Soc. A}, vol. 253, pp. 358--379, 1959.
\href{https://doi.org/10.1098/rspa.1959.0200}{doi:10.1098/rspa.1959.0200}

\end{thebibliography}

\end{document}
